% Options for packages loaded elsewhere
\PassOptionsToPackage{unicode}{hyperref}
\PassOptionsToPackage{hyphens}{url}
%
\documentclass[a4paper
]{ubarticle}
\usepackage{graphicx}
\usepackage{amssymb,amsmath}
    \usepackage{parskip}
  \usepackage{tabu}
  \usepackage{longtable}
\makeatother

\usepackage{hyperref}
\hypersetup{
  pdftitle={Viadi Routing API},
  hidelinks,
  pdfcreator={LaTeX via pandoc}}
\urlstyle{same} % disable monospaced font for URLs
\setlength{\emergencystretch}{3em} % prevent overfull lines
\providecommand{\tightlist}{%
  \setlength{\itemsep}{0pt}\setlength{\parskip}{0pt}}
\setcounter{secnumdepth}{-\maxdimen} % remove section numbering

\title{Viadi Routing API}
\date{\today}
\author{Ubique Innovation AG}

\begin{document}
\begin{titlepage}
	\includegraphics[width=7cm]{ubique-logo.png}
	\hspace{4.3cm}
 {\raggedleft
	 \textbf{Ubique Innovation AG} \\
	\hspace{11.5cm} Niederdorfstrasse 77 \\
	8001 Zürich \\
\vspace{0.3cm}
	\url{www.ubique.ch } 
\par}
	\vspace{3cm}
	{\Huge STARSDK API \par}
	\vspace{1.5cm}
	{\huge Documentation \par}
	\vspace{3cm}
	{ \large \today }
	\end{titlepage}
\thispagestyle{empty}
\clearpage
\tableofcontents
\clearpage
\part{Introduction}
\section{Introduction}
This serves as a technical overview of the backend, with requests and all models. Where possible we provide examples on how the payload could look like.
\part{Web Service}
This part of the documentation deals with the different API-Endpoints. Examples to fields are put into the models section~\ref{sec:Models} to increase readability. Every request lists all possible status codes and the reason for the status code.
\section{ /v1/ }
    \begin{verbatim}
    get /v1/
    \end{verbatim}
hello

\subsection{Responses}
\subsubsection{ 200 Success }
 

    
        \begin{ubresponses}{\textwidth}{|Y|}
        \ubheader{Type}\\
        \hline
             \hyperref[sec:string] { string } \\
 \hline

        \end{ubresponses}
    
\section{ /v1/exposed }
    \begin{verbatim}
    post /v1/exposed
    \end{verbatim}
Enpoint used to publish the SecretKey.

\subsection{ Request Headers }
\begin{ubparam}{\textwidth}{|H|c|Y|}
\ubheader{Field} & \ubheader{Type} & \ubheader{Description}\\
\hline
\ubheader{ User-Agent }   \textcolor{red}{\emph{*}}  &  string  & App Identifier (PackageName/BundleIdentifier) + App-Version + OS (Android/iOS) + OS-Version
 \\
\hline
\end{ubparam}

\subsection{ Response Body }
\begin{ubparam}{\textwidth}{|H|c|Y|}
\ubheader{Field} & \ubheader{Type} & \ubheader{Description}\\
\hline
\ubheader{  }   \textcolor{red}{\emph{*}}  &  ExposeeRequest  & The ExposeeRequest contains the SecretKey from the guessed infection date, the infection date itself, and some authentication data to verify the test result
 \\
\hline
\end{ubparam}
\subsection{Responses}
\subsubsection{ 200 Success }
Returns OK if successful
 

    
        \begin{ubresponses}{\textwidth}{|Y|}
        \ubheader{Type}\\
        \hline
             \hyperref[sec:string] { string } \\
 \hline

        \end{ubresponses}
    
\subsubsection{ 400 Bad Request }
Multiple possible errors: 
\begin{itemize}
\item Key is not base64
\item No AuthData given
\end{itemize} 


\section{ /v1/exposed/{dayDateStr} }
    \begin{verbatim}
    get /v1/exposed/{dayDateStr}
    \end{verbatim}


\subsection{ Path Parameters }
\begin{ubparam}{\textwidth}{|H|c|Y|}
\ubheader{Field} & \ubheader{Type} & \ubheader{Description}\\
\hline
\ubheader{ dayDateStr }   \textcolor{red}{\emph{*}}  &  string  & The date for which we want to get the SecretKey.
 \\
\hline
\end{ubparam}
\subsection{Responses}
\subsubsection{ 200 Success }
Returns ExposedOverview, which includes all secretkeys which were published on \emph{dayDateStr}.
 

    
        \begin{ubresponses}{\textwidth}{|Y|}
        \ubheader{Type}\\
        \hline
             \hyperref[sec:ExposedOverview] { ExposedOverview } \\
 \hline

        \end{ubresponses}
    
\subsubsection{ 500 Internal Server Error }
If the date time string is not parseable.
 



\part{Models}
All Models, which are used by the Enpoints are described here. For every field we give examples, to give an overview of what the backend expects.
\label{sec:Models}
\subsection{ AuthDataMethod }
\label{sec:AuthDataMethod}

\begin{itemize}
    \item NONE
    \item REDEEM\_CODE
\end{itemize}

\subsection{ ExposedOverview }
\label{sec:ExposedOverview}
\begin{ubresponses}{\textwidth}{|H|c|Y|p{2cm}|}
\ubheader{Field} & \ubheader{Type}  &\ubheader{Description}& \ubheader{Example}\\
\hline
 \ubheader{ exposed }  & \hyperref[sec:Exposee]{ Exposee[] }   & A list of all SecretKeys
 &  \\
\hline

\end{ubresponses}

\subsection{ Exposee }
\label{sec:Exposee}
\begin{ubresponses}{\textwidth}{|H|c|Y|p{2cm}|}
\ubheader{Field} & \ubheader{Type}  &\ubheader{Description}& \ubheader{Example}\\
\hline
 \ubheader{ Id }  & \hyperref[sec:integer]{ integer }   &  &  \\
\hline
 \ubheader{ key }  \textcolor{red}{\emph{*}}  & \hyperref[sec:string]{ string }   & The SecretKey of a exposed.
 & AAAA \\
\hline

\end{ubresponses}

\subsection{ ExposeeAuthData }
\label{sec:ExposeeAuthData}
\begin{ubresponses}{\textwidth}{|H|c|Y|p{2cm}|}
\ubheader{Field} & \ubheader{Type}  &\ubheader{Description}& \ubheader{Example}\\
\hline
 \ubheader{ value }  & \hyperref[sec:string]{ string }   & Authentication data used to verify the test result
 &  \\
\hline
 \ubheader{ method }  & \hyperref[sec:AuthDataMethod]{ AuthDataMethod }   & Specify the method, which should be used to verify data
 &  \\
\hline

\end{ubresponses}

\subsection{ ExposeeRequest }
\label{sec:ExposeeRequest}
\begin{ubresponses}{\textwidth}{|H|c|Y|p{2cm}|}
\ubheader{Field} & \ubheader{Type}  &\ubheader{Description}& \ubheader{Example}\\
\hline
 \ubheader{ key }  \textcolor{red}{\emph{*}}  & \hyperref[sec:string]{ string }   & The SecretKey used to generate EphID
 &  \\
\hline
 \ubheader{ authData }  \textcolor{red}{\emph{*}}  & \hyperref[sec:ExposeeAuthData]{ ExposeeAuthData }   & AuthenticationData provided by the health institutes to verify the test results
 &  \\
\hline

\end{ubresponses}



\end{document}
